In this section, we discuss two representative examples illustrating the
challenges of analyzing the precise
\emph{reachability-bound} on
every location.
% We also give the technical overview of our algorithm.
%
\subsection{Multiple-Path Loop}
\label{sec:overview-multiplepath}
\input{chapters/reachability/examples/whileTwoCounters-overview}
Figure~\ref{fig:whileTwoCounters-overview}(a) shows an example of a two paths loops
with different \emph{reachability-bounds} on the locations in different paths.
This example is adopted from the example in~\cite{GulwaniZ10}, which
is a skeleton code from the .Net base-class library.


In this example, given $n \geq m$,
the expected \emph{reachability-bound}s for locations $4$ and $5$ are both $m \times \lfloor\frac{n}{m}\rfloor$,
for location $2$ and $3$ are $(m + 1) \times \lfloor\frac{n}{m}\rfloor + 1$, 
and $1$ for locations $0, 1$ and $\lex$. 
\highlight{Notice here, though within the same loop $L_2$, the bounds for locations $4$ and $5$ on the first branch, and $6$ on the second branch are different.}
State-of-art \emph{reachability-bound} analysis~\cite{GulwaniZ10}
gives the same \emph{reachability-bound}, $n + \lfloor\frac{n}{m}\rfloor$ for all the locations within the loop $L_2$, which is tight w.r.t. $L_2$'s iteration times but not for different locations inside $L_2$ without considering multiple paths.
Among works on program complexity, cost, and loop bound analysis, the algorithm in~\cite{GulwaniJK09} computes a tight bound on the loop iteration times but not precise \emph{reachability-bound}s on different locations.
Though we can use this loop iteration bound as the \emph{reachability-bound}s for locations $1$ and $2$,
the bounds for locations $4, 5$ and $6$ are still unclear.


\paragraph{Path Reachability-bound}
To compute the bounds for locations on different paths of a loop, we compute a new quantity -- \emph{path reachability-bound}.
% which is the first key idea of this path-sensitive \emph{reachability-bound} analysis algorithm. 
This quantity bounds the evaluation times of each \emph{simple transition path} during the program execution, where a \emph{simple transition path} is loop-free and interleaving-free.

For this example, we first
generate the abstract transition graph for the program as in Figure~\ref{fig:whileTwoCounters-overview}(b), using the difference constraint-based program abstraction method.
Then we refine the loop in this program by checking the 
execution orders of each simple transition path and
generates a refined program $\rprog$ as below,
% 
\[ \tpath_0 ; 
 \rpchoose{
    \rprepeat(\rprepeat(\tpath_1); \tpath_2), 
   \rprepeat(\tpath_1)}; \tpath_3.
\]
In this refined program,
each $\tpath_i$ is a \emph{simple transition path} we computed in a pre-procedure, which is loop-free and not interleaving with the other paths, $\tpath_j, j \neq i$ as in Figure~\ref{fig:whileTwoCounters-overview}(c).
$\rprepeat(\tpath_1)$ represents that execution of this program contains iteration on the path $\tpath_1$.
$\rprepeat(\tpath_1); \tpath_2$ denotes that after the iterations of $\tpath_1$ is done,
$\tpath_2$ will execute right after.
$\rpchoose{
    \rprepeat(\rprepeat(\tpath_1); \tpath_2), 
   \rprepeat(\tpath_1)}$ means that there are two possible execution orders for the original loop $L_2$.
The execution orders of these \emph{simple transition path}s are explicated in the refined program.
% Every path will not interleave with the others. 
Then we compute the \emph{path reachability-bound} for every $\tpath_i$,
$\inoutB(\rprog, \tpath_i)$ during the execution of $\rprog$ through the ranking function estimation in Section~\ref{sec:psRB-rank} and the algorithm in Section~\ref{sec:psRB-pathlocalrb} to compute the result as
% which is a bound on the reachability time of $\tpath$ during the execution of $\rprog$.
% The \emph{path reachability-bound}s for the four simple transition paths in this example are
$\inoutB(\rprog, \tpath_1) = \max\{m, m \times \frac{n}{m} \}$,
$\inoutB(\rprog, \tpath_2) = \frac{n}{m}$,
and $\inoutB(\rprog, \tpath_0) = \inoutB(\rprog, \tpath_3) = 1$.
% \\
% and another \emph{loop reachability-bound}
% to compute the final \emph{reachability-bound} for each location.
Since there isn't nested loop in this example, we simply sum up $\inoutB(\rprog, \tpath)$ over the $\tpath$ where a certain location shows up
and as the \emph{reachability-bound} of this location.
Then we get the precise \emph{reachability-bound} for every location in program $\kw{twoPathsWhile}$ as
$\psRB(0) = \psRB(1) = \psRB(\lex) = 1$,
$\psRB(4) = \psRB(5) = \max\{m, m \times \lfloor\frac{n}{m}\rfloor\}$,
$\psRB(3) = \psRB(2) = \max\{m, m \times \lfloor\frac{n}{m}\rfloor\} + \lfloor\frac{n}{m}\rfloor + 1 $,
and $\psRB(6) = \lfloor\frac{n}{m}\rfloor$.
%
\subsection{Nested Loops with Related Iterator}
\label{sec:overview-nestedwhile}
However, when there exists a nested loop, computing the \emph{reachability-bound} for each location encounters another challenge.
% The \emph{path reachability-bound} is precise for each path w.r.t. the innermost loop but not the outer loops.
Figure~\ref{fig:threeWhile-overview}(a) shows an example of the nested loops with related 
iterators.
This example is adopted from the example in~\cite{GulwaniJK09}, which is common in product code.

In line 8, $i$ is reset by $w$, and $w$ is reset by $j$ in line 5. So the
while $L_6$ is only executed in the first iteration of while loop $L_1$ and $L_3$.
% \\
The while loop $L_3$ at line 3 is executed only in 
the first $m - N$ iterations of the 
$L_1$ because $j$ is reset by $i$ in line 2.
% \\
So the total iterations of all three loops are
$n + m^2 - m \times N$,
and the precise \emph{reachability-bound} for location $7$ inside the $L_6$ is $N$,
for locations $4, 5$ and $8$ between the $L_3$ and $L_6$ are $(n-N) \times (m - N)$,
and $n - N$ for locations $2$ and $9$.
% \\
Notice here, the \emph{reachability-bounds} for the locations inside the loop $L_6$ is 
the same as its innermost loop iteration bound.
% , as well as our \emph{path reachability-bound}.
However, for the locations between $L_3$ and $L_6$,
the \emph{reachability-bounds} are the multiplication of the inner and outer loop iteration bounds.
To the best of our knowledge, the loop bound analysis method in \cite{GulwaniJK09} can only give a loose upper bound $n + (m \times n) + N$ for the entire loop complexity, and 
the difference constraint-based algorithm in \cite{SinnZV17} is able to
compute a better but still loose bound, $n + m^2 - m \times N$ on total iteration times.
None of them can give the precise \emph{reachability-bound} for every location in these nested loops,
which is non-trivial to compute even though knowing the loop bound,
especially for the locations similar to $7$ in $\kw{threeNestedWhile}$.

In order to precisely compute how many times the location $7$ is reached, 
the local \emph{path reachability-bound} we discussed in the first example is not precise enough.
We need to know
the numbers of iterations of the outer loop $L_3$ and $L_1$ such that,
during these iterations, the inner loop $L_6$ is ``reached''. 
We call the upper bound on this number, the \emph{loop reachability-bound} of the location within loop $L_6$ w.r.t the outer loops $L_3$ and $L_1$.
Then by multiplying the loop iteration bound of the $L_6$ with its \emph{loop reachability-bound} w.r.t the $L_3$ and $L_1$, we can compute the precise
\emph{reachability-bounds} for location $7$.

The method in \cite{GulwaniJK09} based on the \emph{progress invariant} computes a similar quantity, but it only computes
the
bound on iteration numbers
of the inner loop $L_6$ in each iteration of $L_3$ and $L_1$, which are both $N$.
So they have to over-approximate the reachability-bound for locations inside $L_6$ with the
overall program complexity by multiplication, i.e., $n + m^2 - m \times N$.
% \\
The difference constraint-based algorithm in~\cite{SinnZV17} through ranking function estimation
is only able to
compute the combined loop bound and the local bound of each loop
separately as well.
% We are still unable to know the precise \emph{reachability-bound} for the locations in the innermost loop.

\paragraph{Loop reachability-bound.}
In this sense, we introduce the \emph{loop reachability-bound}
% , $\lpchB(\rprog, \tpath)$.
% is our second key idea combining two lines of works.
% It has high accuracy and efficiency by using the estimated ranking function based on the \emph{amortized complexity analysis} methodology over the refined loop paths.
For each \emph{simple transition path}  w.r.t each of its outer loop,
% in which $\tpath$ is nested,
\emph{loop reachability-bound}
is a bound on the iterations of
the outer loop
% $L:\rprog$ w.r.t. the innermost loop where $\tpath$ is enclosed,
such that during these outer loop iterations, the inner loop is ``reached''. 
This is distinguished from the traditional methods, which only estimate the bound on the inner loop's iteration number
in one iteration of the outside loop.


Similar to the $\kw{twoPathsWhile}$ example, we generate its abstract transition graph as in Figure~\ref{fig:threeWhile-overview}(a),
and compute the refined program,
$\rprog = \tpath_0; 1: \rprepeat(\tpath_1;$ 
$3: {\rprepeat(\tpath_2; 6 : {\rprepeat(\tpath_3)}; \tpath_4)}; \tpath_5);$ 
$\tpath_6$,
where the $\tpath_0, \ldots$ are shown in the middle part of Figure~\ref{fig:threeWhile-overview}(b).
We use $\rprog_1$ and $\rprog_3$ denote the body of the loop $L_1$ and $L_3$ respectively as in the bottom part of Figure~\ref{fig:threeWhile-overview}(b).
% to denote ${\rprepeat(\tpath_1; 3: {\rprepeat(\tpath_2; 6 : {\rprepeat(\tpath_3)}; \tpath_4)}; \tpath_5)}$
% and $\rprog_3 = {\rprepeat(\tpath_2; 6 : {\rprepeat(\tpath_3)}; \tpath_4)}$
In the first step, we still compute the \emph{path reachability-bound} locally for each $\tpath_i$ but only w.r.t. the innermost loop in which $\tpath_i$ is nested.
Then differently from $\kw{twoPathsWhile}$,
we compute \emph{loop reachability-bound}, $\lpchB(\rprog_j, \tpath_i)$ for each $\tpath_i$ w.r.t. each of its outer loop by the algorithm in Section~\ref{sec:psRB-looprb}.
For $\tpath_3$ we compute
$\lpchB(1: \rprog_1, \tpath_3) = 1$ and
$\lpchB(3: \rprog_3, \tpath_3) = 1$.
Both are tight because loop $L_6$ will only be entered once among all iterations of $L_1$ and $L_3$, and in all the rest iterations, the body of loop $L_6$ isn't executed at all.
% So $1$ as \emph{loop reachability-bound} of this path is tight w.r.t. both the loop $L_3$ and $L_1$.
% In the same way, we also compute $\lpchB(3: \rprog_3, \tpath_3) = 1$ precisely.
Then for each $\tpath_i$, the multiplication of its local \emph{path reachability-bound} with all its \emph{loop reachability-bound}s is the accurate \emph{path reachability-bound} globally.
By summing up the reachability-bound of the path where each location shows up,
% as its \emph{reachability-bound} as before.
% and multiply this result by all its \emph{loop reachability-bound}s.
% In this way, 
we compute $N$ as the \emph{reachability-bound} of location $7$, which is tight.
\input{chapters/reachability/examples/threeNestedWhile-overview}